% Generated draft: requires mathematical and source review.
\documentclass[11pt]{article}
\usepackage[margin=25mm]{geometry}
\usepackage{fontspec}
\setmainfont{Latin Modern Roman}
\usepackage{amsmath,amssymb,mathtools,unicode-math}
\setmathfont{Latin Modern Math}
\usepackage{enumitem,needspace,xurl,hyperref}
\setlength{\parindent}{0pt}
\setlength{\parskip}{5pt}
\setlength{\emergencystretch}{3em}
\newcommand{\N}{\mathbb{N}}
\newcommand{\Z}{\mathbb{Z}}
\newcommand{\Q}{\mathbb{Q}}
\newcommand{\R}{\mathbb{R}}
\newcommand{\C}{\mathbb{C}}
\newcommand{\F}{\mathbb{F}}
\newcommand{\A}{\mathbb{A}}
\newcommand{\PP}{\mathbb{P}}
\newcommand{\E}{\mathbb{E}}
\newcommand{\eps}{\varepsilon}
\newcommand{\dd}{\,\mathrm{d}}
\begin{document}
\section*{UnsolvedMath ID: 30001802}
\textbf{Source problem code:} catalogue entry
\subsection{Definitions and mathematical statement}
For terminology only, let \(A\in K^{n\times n}\), where \(K\) is a field, and let \(L\) be an extension field of \(K\). An eigenvalue of \(A\) in \(L\) is a scalar \(\lambda\in L\) for which there exists \(v\in L^n\setminus\{0\}\) satisfying
\[
Av=\lambda v.
\]
Such a vector \(v\) is an eigenvector, and the corresponding eigenspace is \(\ker(A-\lambda I_n)\). These definitions supply vocabulary only; they do not specify the matrix ensemble intended by the source.

The retained source question is: \emph{What do the eigenvalues and eigenvectors of a typical large random matrix look like?}

\textbf{Unresolved formulation.} The supplied statement-verification record identifies this sentence as a rhetorical introduction to a survey of established spectral and eigenvector universality results. It records the statement as \texttt{unrecoverable}: no unresolved conjecture, matrix ensemble, scale, or target theorem is specified, and no separate clean statement is supplied. In particular, the source does not determine the underlying field, the probability distributions of the matrices, structural or moment assumptions, the meaning of ``typical,'' a limiting regime, spectral normalization, eigenvector normalization, observables, or a mode of convergence. Consequently, no precise quantified conjecture can be recovered from this record without additional choices. The registry's supplied \texttt{open} label does not by itself identify an unresolved mathematical proposition.
\subsection{Short English statement}
The record asks broadly about eigenvalues and eigenvectors of large random matrices. Its supplied verification identifies an expository framing question, with no specified ensemble or conjecture. A precise mathematical problem cannot be extracted without additional hypotheses and a target conclusion; the supplied \texttt{open} label requires this qualification.
\subsection{Sources}
\begin{itemize}
\item Local UnsolvedMath registry, \texttt{data/unsolvedmath/problems.json}, permanent ID \(30001802\), title \emph{Expository Random-Matrix Framing Question}, source code \texttt{catalogue entry}, in the Oberwolfach Reports Open Problems collection. The supplied record gives proposed year 2011 and names no individual proposer.
\item Supplied report citation: \emph{Stochastic Analysis} (2011), Oberwolfach Reports, pp.\ 1573--1635, DOI \texttt{10.4171/owr/2011/29}. \url{https://doi.org/10.4171/owr/2011/29}. This is also the supplied literature source labeled ``original Oberwolfach report.'' No narrower passage location is supplied.
\item Supplied literature source labeled ``Erdos--Yau--Yin bulk universality'': \url{https://arxiv.org/abs/1001.3453}.
\item Supplied dated literature curation, checked \texttt{2026-08-21}, labels the record \texttt{open}. Its assessment says that the broad question is not a single well-posed universality claim; it reports bulk and edge laws, rigidity, and eigenvector delocalization for broad Wigner classes, while stating that no general notion of a typical large matrix resolves the extracted question. This is supplied dated curation, not an independently verified literature assessment. The background expressly requests independent verification of time-sensitive status and cited claims.
\item Supplied statement-verification note: the question introduces a survey of established spectral and eigenvector universality results, rather than specifying an unresolved conjecture. The original record is retained for provenance, with statement status \texttt{unrecoverable} and no promoted clean statement.
\end{itemize}
\end{document}
